\chapter*{Abstract}
\addcontentsline{toc}{chapter}{Abstract}

Modern academic research depends on a fragmented stack of tools. Researchers locate papers in one application, manage references in a second, draft documents in a third, edit collaboratively in a fourth, and consult AI assistants in a fifth. Each transition between these tools requires the user to rebuild context, copy artifacts manually, and reconcile inconsistent state. Recent work shows that this fragmentation imposes a measurable cognitive cost on researchers and that general-purpose AI assistants make the problem worse by fabricating bibliographic citations at rates between 18\% and 55\% depending on model and prompt \cite{walters2023fabrication}.

This work presents \sh{}, a unified web-based research workspace that integrates paper discovery, reference management, real-time collaborative writing, and AI assistance grounded in the user's library. The system is implemented as a three-tier architecture: a FastAPI backend in Python (approximately 60{,}000 lines across 189 modules), a React and TypeScript frontend (approximately 58{,}000 lines across 198 components), and a Hocuspocus collaboration server backed by Y.js conflict-free replicated data types. Paper discovery aggregates results across nine open scientific APIs (Semantic Scholar, OpenAlex, arXiv, CrossRef, PubMed, CORE, Europe PMC, ScienceDirect, and Google Scholar) with deduplication based on DOI and normalized title matching. The AI subsystem uses a tool-orchestrated agent loop with thirty registered tools spanning library access, search, paper analysis, and content creation, allowing the assistant to ground every citation in the user's curated reference library.

The system is evaluated through three studies. First, a citation faithfulness benchmark of 30 cited-writing prompts against four frontier 2026 LLMs (GPT-5 Mini, Claude Haiku 4.5, Gemini 3 Flash, and DeepSeek R1) in two conditions (ungrounded baseline and library-grounded) measures citation existence and library faithfulness with external verification through the Semantic Scholar and OpenAlex APIs. Second, a system-performance evaluation against the deployed backend measures end-to-end latencies for paper-search aggregation, AI orchestrator first-token time, and parallel-versus-sequential discovery. Third, a comparative workflow analysis estimates, via the Keystroke-Level Model with empirically measured per-source network round-trips, the reduction in user actions required to complete five canonical research tasks against the conventional three-tool stack of Overleaf, Zotero, and ChatGPT.

The contributions of this work are six, summarized in chapter~\ref{ch:conclusion} as C1--C6. The headline empirical results: over a 30-prompt benchmark per cell (240 model calls in total, 1{,}135 verified citations), the tool-orchestrated AI agent reduces mean raw-model citation fabrication from 76.3\% (ungrounded) to 14.4\% (grounded) across the four 2026 frontier models, and the production citation filter --- which operates only on LaTeX \texttt{\textbackslash cite\{key\}} commands and rejects any key not derivable from a curated library entry --- drives the user-visible fabrication rate to 0.0\% across all four models (329 library-key citations kept, 0 fabricated). The system has been deployed at \url{https://scholarhub.space} during the development of this work.

\vspace{0.5cm}
\noindent\textbf{Keywords:} Collaborative editing, conflict-free replicated data types, retrieval-augmented generation, academic writing tools, large language model citation, paper discovery, multi-source aggregation.

\clearpage
